% =============================================================================
\chapter*{Resumo}
\addcontentsline{toc}{chapter}{Resumo}

% Reduz o espaço vertical após o título e amplia localmente a área útil.
\vspace{-1.2cm}
\enlargethispage{2\baselineskip}

A contaminação ambiental por compostos prioritários exige a articulação de
informações sobre estruturas químicas, funções biológicas, vias de degradação,
toxicidade e referências institucionais. Essas informações permanecem
distribuídas entre bases com unidades de observação e critérios de curadoria
distintos, o que dificulta a interpretação de perfis funcionais derivados de
dados ômicos. O BioRemPP constitui um ecossistema computacional orientado por
compostos que integra esse conteúdo e o disponibiliza por diferentes formas de
consulta e análise. O ecossistema reúne quatro componentes: o banco de dados
curado BioRemPP, o BioRemPP Database Explorer, o BioRemPP Web Service e uma
interface de linha de comando. O banco integra a tabela central BioRemPP, o
HADEG, o subconjunto KEGG Degradation e o ToxCSM por meio de relações entre
grupos de ortologia KEGG (KOs), compostos, números EC, reações, vias,
descritores de atividade enzimática, referências institucionais e predições
toxicológicas. Um workflow Snakemake automatiza sua construção, enquanto
verificações ancoradas nas relações do KEGG e validações estruturais e
regressivas com Great Expectations avaliam a consistência dos artefatos
gerados. O Database Explorer oferece cinco módulos de exploração e oito
análises guiadas, acompanhadas por consultas SQL versionadas. O Web Service
recebe conjuntos de KOs organizados por amostra e reúne 56 casos de uso
distribuídos em oito módulos, que abrangem análises descritivas,
multivariadas, relacionais, toxicológicas e de complementaridade funcional. A
interface de linha de comando disponibiliza as mesmas operações para execução
programática e incorporação a pipelines automatizados. A versão 1.1.0 do banco
reúne 384 compostos distribuídos em 12 classes químicas, 1.543 KOs, 994 números
EC, 1.939 identificadores de reação KEGG, 1.517 símbolos gênicos e 206
descritores de atividade enzimática. Nove códigos de referências institucionais
delimitam o recorte de compostos prioritários, e 370 compostos apresentam
predições para 31 endpoints toxicológicos. Os componentes preservam os
identificadores KO e CPD, registram versões e parâmetros de análise e fornecem
tabelas, visualizações e artefatos exportáveis. O BioRemPP apoia a triagem
\textit{in silico}, a comparação de perfis funcionais e a seleção de hipóteses
para estudos posteriores. As associações baseadas na presença de KOs e os
valores toxicológicos preditos representam potencial funcional e contexto
computacional, mas não substituem validações transcriptômicas, proteômicas,
metabolômicas, bioquímicas ou ambientais.

\vspace{0.35em}
\noindent\textbf{Palavras-chave:} bioinformática; biorremediação; genômica
funcional; integração de dados ômicos; compostos prioritários; ecossistema de
software.